\section{Results}
\label{sec:results}

Results are organized around five claims. The headline result:
TTA with adaptive deferral reduces segmentation error by approximately
80\% at a 25\% deferral rate (Table~\ref{tab:main_results}).
Figure~\ref{fig:summary_scatter} plots error reduction against deferral
rate for all method--policy configurations, making the interaction between
these two design choices visible at a glance.

\begin{figure}[H]
\centering
\includegraphics[width=0.85\textwidth]{summary_scatter.png}
\caption{The deferral design space. Each point is one uncertainty method
paired with one deferral policy. The upper-left corner (high error reduction,
low deferral rate) is ideal. TTA+adaptive provides the largest absolute
error reduction, while TTA+confidence-aware is the strongest low-budget
operating point.}
\label{fig:summary_scatter}
\end{figure}

\begin{table}[H]
\centering
\caption{Main results on the DRIVE test set (20 images). Best values in bold.
Runtime is per image on a single GPU. ERR = error reduction ratio.
``Def.\ \%'' = percentage of pixels deferred.}
\label{tab:main_results}
\small
\setlength{\tabcolsep}{4.6pt}
\resizebox{\textwidth}{!}{%
\begin{tabular}{lcccccccc}
\toprule
\textbf{Method} & \textbf{Dice} & \textbf{AUC} & \textbf{IoU} &
\textbf{ECE} & \textbf{Unc-AUROC} & \textbf{ERR} &
\textbf{Def.\ \%} & \textbf{Time (s)} \\
\midrule
Deterministic       & 0.736 & 0.944 & --- & 0.245 & 0.500 & --- & --- & 0.08 \\
MC Drop.\ (global)  & 0.763 & 0.965 & 0.617 & 0.035 & 0.722 & 26.5\% & 33.4\% & 2.65 \\
MC Drop.\ (adaptive) & 0.763 & 0.965 & 0.617 & 0.035 & 0.722 & 42.4\% & 25.0\% & 2.65 \\
MC Drop.\ (conf-aware) & 0.763 & 0.965 & 0.617 & 0.035 & 0.722 & 48.8\% & 10.9\% & 2.65 \\
TTA (global)        & 0.768 & 0.969 & 0.624 & \textbf{0.035} & \textbf{0.881} & 55.1\% & 13.2\% & \textbf{0.86} \\
TTA (adaptive)      & \textbf{0.768} & \textbf{0.969} & \textbf{0.624} & \textbf{0.035} & \textbf{0.881} & \textbf{79.5\%} & 25.0\% & \textbf{0.86} \\
TTA (conf-aware)    & 0.768 & 0.969 & 0.624 & \textbf{0.035} & \textbf{0.881} & 55.1\% & 12.1\% & \textbf{0.86} \\
\midrule
Ensemble (N=5)      & 0.771 & 0.947 & 0.627 & 0.031 & 0.833 & --- & --- & 3.90 \\
\bottomrule
\end{tabular}%
}
\end{table}

\subsection{TTA produces better uncertainty than MC Dropout}
\label{sec:results_uncertainty}

\textbf{Claim:} TTA uncertainty discriminates correct from incorrect
predictions substantially better than MC Dropout uncertainty, while
requiring less computation.

TTA achieves an Unc-AUROC of 0.881, compared to 0.722 for MC Dropout
(Table~\ref{tab:main_results}). This is a large gap: TTA correctly ranks
88.1\% of error/non-error pixel pairs by uncertainty, versus 72.2\% for
MC Dropout. The difference is consistent across images (bootstrap 95\%
CI for the external-dataset MC Dropout analysis remains well above chance,
and the unified comparison bundle shows the same ranking gap clearly).

TTA also produces higher AUCC for both Dice (0.846 vs.\ 0.792) and AUC
(0.888 vs.\ 0.884), meaning that progressively removing the most uncertain
pixels improves performance faster with TTA uncertainty.

In terms of computation, TTA requires 6 forward passes with deterministic
augmentations ($K = 6$), while MC Dropout requires 30 stochastic passes
($T = 30$). TTA inference takes about 0.86\,s per image in the unified
comparison bundle, while MC Dropout takes about 2.65\,s, a 3.1$\times$ speedup.

\paragraph{Interpretation.} The quality gap between TTA and MC Dropout
uncertainty has a structural explanation. In retinal vessel segmentation,
the most common errors occur at vessel boundaries and thin vessel tips,
where the segmentation mask is sensitive to small input changes. TTA
directly tests this sensitivity by applying geometric transforms; if a
boundary pixel flips under rotation, TTA flags it. MC Dropout instead
samples weight perturbations, which produce more diffuse uncertainty that
does not concentrate as sharply on error-prone structures.

\begin{figure}[H]
\centering
\includegraphics[width=\textwidth]{uncertainty_maps.png}
\caption{Comparison of uncertainty maps. TTA uncertainty concentrates on
boundary and thin-vessel regions where errors are most frequent, while
MC Dropout produces more spatially diffuse uncertainty.}
\label{fig:uncertainty_maps}
\end{figure}

\subsection{Adaptive maximizes error removal, confidence-aware is best at low budget}
\label{sec:results_deferral}

\textbf{Claim:} Adaptive deferral gives the largest absolute error removal,
while confidence-aware deferral is the most efficient low-budget policy.

Table~\ref{tab:deferral_comparison} breaks down the deferral results. For MC
Dropout, global thresholding defers 33.4\% of pixels and reduces error by
26.5\%. Adaptive deferral improves that to 42.4\% at the same 25\% review
budget used by the adaptive TTA policy, and confidence-aware deferral reaches
48.8\% while deferring only 10.9\% of pixels.

For TTA, adaptive deferral yields the headline result: 79.5\% error reduction
at 25.0\% deferred pixels. The confidence-aware policy reaches essentially the
same error reduction as the global TTA policy, but with a smaller review budget
(12.1\% vs.\ 13.2\% deferred).

\begin{table}[H]
\centering
\caption{Deferral strategy comparison on DRIVE test set. Error rates are
pixel-wise misclassification rates. Adaptive deferral maximizes absolute error
removal, while confidence-aware deferral gives the strongest reduction per
deferred percentage point at low review budgets.}
\label{tab:deferral_comparison}
\small
\setlength{\tabcolsep}{4.6pt}
\resizebox{\textwidth}{!}{%
\begin{tabular}{llccccc}
\toprule
\textbf{Uncertainty} & \textbf{Deferral} & \textbf{Error before} &
\textbf{Error after} & \textbf{ERR} & \textbf{Def.\ \%} &
\textbf{ERR/Def.\ \%} \\
\midrule
MC Dropout & Global     & 0.0658 & 0.0484 & 26.5\% & 33.4\% & 0.79 \\
MC Dropout & Adaptive   & 0.0658 & 0.0379 & 42.4\% & 25.0\% & 1.70 \\
MC Dropout & Conf-aware & 0.0658 & 0.0337 & 48.8\% & 10.9\% & 4.49 \\
\midrule
TTA        & Global     & 0.0640 & 0.0287 & 55.1\% & 13.2\% & 4.19 \\
TTA        & Adaptive   & 0.0640 & 0.0131 & 79.5\% & 25.0\% & 3.18 \\
TTA        & Conf-aware & 0.0640 & 0.0288 & 55.1\% & 12.1\% & 4.56 \\
\bottomrule
\end{tabular}%
}
\end{table}

The rightmost column in Table~\ref{tab:deferral_comparison}, ERR per
percentage point deferred (ERR/Def.\ \%), quantifies deferral efficiency.
For MC Dropout, confidence-aware deferral is dramatically more efficient than
the global policy (4.49 vs.\ 0.79). For TTA, both global and confidence-aware
policies are efficient, with confidence-aware remaining slightly better
(4.56 vs.\ 4.19) at a lower deferral rate.

\paragraph{Interpretation.} Global thresholding wastes deferral budget on
pixels that are uncertain but correctly predicted. As formalized in
Section~\ref{sec:deferral}, the confidence-aware score $s = u \cdot (1 - c)$
suppresses these false alarms. Vessel interiors are the clearest example: MC
Dropout assigns moderate uncertainty there, but predictions sit well above
0.5, so $c$ is high and $s$ stays low.

\begin{figure}[H]
\centering
\includegraphics[width=\textwidth]{error_reduction_bars.png}
\caption{Error rates before and after deferral. TTA+adaptive achieves
80\% error reduction, while TTA+conf-aware achieves 55\% error reduction
at about half the adaptive review budget. The difference between choosing the
right uncertainty method and the right policy is large enough to reshape the
practical review workflow.}
\label{fig:error_reduction_bars}
\end{figure}

\begin{figure}[H]
\centering
\includegraphics[width=\textwidth]{deferral_curves.png}
\caption{Coverage vs.\ Dice for all method--policy combinations. TTA
curves dominate MC Dropout curves at every coverage level. Among deferral
policies, confidence-aware reaches clinical-threshold Dice (0.82) at the
highest coverage, meaning the fewest pixels need review.}
\label{fig:deferral_curves}
\end{figure}

\begin{figure}[H]
\centering
\includegraphics[width=\textwidth]{unc_vs_error.png}
\caption{Error rate by uncertainty decile. TTA uncertainty sorts pixels
into more informative risk strata: the highest-uncertainty decile contains
45\% errors versus 22\% for MC Dropout. This sharper separation
is why TTA enables more efficient deferral.}
\label{fig:unc_vs_error}
\end{figure}

\subsection{Risk-coverage analysis}
\label{sec:results_risk_coverage}

\textbf{Claim:} Selective prediction using TTA uncertainty produces a
steeper risk-coverage curve than MC Dropout, and achieves clinically
acceptable Dice ($\geq 0.82$) at 90\% coverage.

Figure~\ref{fig:risk_coverage} (described below) shows Dice as a function
of coverage for both uncertainty methods. At full coverage (no deferral),
TTA Dice is 0.768 and MC Dropout is 0.764. As coverage decreases (more
pixels deferred), both improve, but TTA improves faster.

At 90\% coverage (deferring 10\% of pixels):
\begin{itemize}[leftmargin=*]
  \item MC Dropout: Dice = 0.778
  \item TTA: Dice = 0.849
\end{itemize}

The clinical threshold for retinal vessel segmentation is typically taken as
Dice $\geq 0.82$ \citep{fraz2012ensemble}. In the unified comparison bundle,
TTA reaches this threshold at roughly 92.7\% coverage, while MC Dropout does
not reach it until roughly 71.8\% coverage. This gap is practically
significant: the TTA operating point retains most pixels while clearing the
clinical threshold, whereas MC Dropout needs much heavier abstention.

The AUCC values quantify this difference:
\begin{itemize}[leftmargin=*]
  \item TTA AUCC (Dice): 0.846 vs.\ MC Dropout: 0.801
  \item TTA AUCC (AUC): 0.888 vs.\ MC Dropout: 0.882
\end{itemize}

\begin{figure}[H]
\centering
\includegraphics[width=\textwidth]{risk_coverage.png}
\caption{Risk-coverage curves for MC Dropout and TTA. TTA reaches
clinically acceptable Dice (0.82) at substantially higher coverage than
MC Dropout in our experiments.}
\label{fig:risk_coverage}
\end{figure}

\paragraph{Operating point analysis.} Table~\ref{tab:operating_points}
shows three clinically motivated operating points for each method.

\begin{table}[H]
\centering
\caption{Operating points for selective prediction on DRIVE.
``Balanced F1'' maximizes the deferral F1 score. ``High precision''
minimizes false deferral (wasted review effort). ``Clinical target''
achieves Dice $\geq$ 0.82 with minimum deferral.}
\label{tab:operating_points}
\small
\setlength{\tabcolsep}{4.6pt}
\resizebox{\textwidth}{!}{%
\begin{tabular}{llcccc}
\toprule
\textbf{Method} & \textbf{Operating point} & \textbf{Coverage} &
\textbf{Dice (acc)} & \textbf{Def. precision} & \textbf{Def. recall} \\
\midrule
MC Dropout & Balanced F1     & 71.9\% & 0.855 & 0.318 & 0.803 \\
MC Dropout & High precision  & 77.3\% & 0.834 & --- & --- \\
MC Dropout & Clinical target & 80.0\% & 0.820 & --- & --- \\
\midrule
TTA        & Balanced F1     & 80.0\% & 0.838 & --- & --- \\
TTA        & Clinical target & 90.0\% & 0.825 & --- & --- \\
TTA        & Aggressive      & 75.0\% & 0.867 & --- & --- \\
\bottomrule
\end{tabular}%
}
\end{table}

\subsection{Calibration and deferral quality are decoupled}
\label{sec:results_calibration}

\textbf{Claim:} Better calibration does not produce better deferral.
In our experiments, temperature scaling worsens the reported ECE
values and leaves the achieved deferral operating points unchanged. This
illustrates that calibration metrics and selective prediction quality should be
evaluated separately.

Table~\ref{tab:calibration} compares calibrated and uncalibrated versions
of each uncertainty method.

\begin{table}[H]
\centering
\caption{Effect of temperature scaling on calibration and deferral quality.
$T$ is the learned temperature. Unc-AUROC measures how well uncertainty
predicts segmentation error (higher is better). In the main comparison bundle,
temperature scaling worsens the reported ECE values and leaves the
uncertainty-ranking metric essentially unchanged.}
\label{tab:calibration}
\small
\setlength{\tabcolsep}{4.6pt}
\resizebox{\textwidth}{!}{%
\begin{tabular}{lcccccc}
\toprule
\textbf{Method} & $T$ & \textbf{ECE (before)} & \textbf{ECE (after)} &
$\Delta$\textbf{ECE} & \textbf{Unc-AUROC (before)} &
\textbf{Unc-AUROC (after)} \\
\midrule
MC Dropout & 1.20 & 0.0350 & 0.0430 & +0.0080 & 0.722 & 0.722 \\
TTA        & 1.35 & 0.0345 & 0.0393 & +0.0047 & 0.881 & 0.881 \\
\bottomrule
\end{tabular}%
}
\end{table}

For MC Dropout, temperature scaling increases the reported ECE from 0.0350 to
0.0430 while leaving the uncertainty-ranking metric unchanged at 0.722.
The deferral operating point we tracked is unchanged as well.

For TTA, the situation is similar: ECE increases from 0.0345 to 0.0393 after
calibration, while Unc-AUROC remains at 0.881. The temperature-scaled run does
not improve the downstream error-reduction operating point.

\paragraph{Why this happens.} Both models learn $T > 1$, confirming mild
overconfidence. But ECE is already below 0.04 before calibration in our
experiments, leaving little room for improvement. The learned
temperatures mostly rescale probabilities without changing the underlying
uncertainty maps used for ranking and deferral, so the ranking metrics stay
flat while the reported calibration values can move in an unfavorable
direction.

The lesson is concrete: if the end use is deferral, evaluate temperature
scaling by Unc-AUROC and achieved operating points, not ECE alone.

\begin{figure}[H]
\centering
\includegraphics[width=0.49\textwidth]{reliability_mc.png}
\includegraphics[width=0.49\textwidth]{reliability_tta.png}
\caption{Reliability diagrams for MC Dropout and TTA, with and without
temperature scaling. The pre-calibration models are already close to
the diagonal, leaving little room for temperature scaling to improve.}
\label{fig:reliability}
\end{figure}

\subsection{Cross-dataset evaluation}
\label{sec:results_cross_dataset}

\textbf{Claim:} Uncertainty quality transfers more robustly than segmentation
accuracy under domain shift.

Table~\ref{tab:cross_dataset} shows results when the DRIVE-trained model
is evaluated on STARE and CHASE\_DB1 without any fine-tuning.

\begin{table}[H]
\centering
\caption{Cross-dataset evaluation (zero-shot transfer from DRIVE).
Segmentation quality drops by 4--5\% Dice, but uncertainty quality
(Unc-AUROC) remains above 0.83 on both external datasets, within
5\% of its DRIVE value.}
\label{tab:cross_dataset}
\small
\setlength{\tabcolsep}{4.6pt}
\resizebox{\textwidth}{!}{%
\begin{tabular}{lccccc}
\toprule
\textbf{Dataset} & \textbf{Dice} & \textbf{AUC} & \textbf{ECE} &
\textbf{Unc-AUROC} & \textbf{AUCC (Dice)} \\
\midrule
DRIVE (in-domain) & 0.791 & 0.960 & 0.043 & 0.850 & 0.793 \\
STARE             & 0.727 & 0.949 & 0.038 & 0.846 & 0.773 \\
CHASE\_DB1        & 0.726 & 0.952 & 0.035 & 0.835 & 0.766 \\
\bottomrule
\end{tabular}%
}
\end{table}

Dice drops from 0.791 (DRIVE) to 0.727 (STARE) and 0.726 (CHASE\_DB1), a
roughly 6.5-point decline. Unc-AUROC, however, remains strong: 0.846 on
STARE and 0.835 on CHASE\_DB1, both close to the in-domain value. This
counterintuitive
result makes sense once you consider that on out-of-distribution data, the
model makes more errors, and those errors tend to occur in regions where the
model is also more uncertain (unfamiliar appearance patterns, different
vessel calibers). The correlation between uncertainty and error therefore
remains useful even as segmentation quality drops.

Put differently, on the images where deferral matters most (unfamiliar
data, lower accuracy), the uncertainty signal is at its most reliable.
This property is directly useful: a system deployed on a new scanner or
patient population can lean on deferral even as segmentation degrades.

\begin{figure}[H]
\centering
\includegraphics[width=\textwidth]{uncertainty_maps.png}
\caption{Qualitative results across three datasets. Errors concentrate at
vessel boundaries (column d). TTA uncertainty tracks these regions (column e),
and the deferral policy routes them for review (column f) while retaining
the large majority of correctly segmented pixels.}
\label{fig:qualitative}
\end{figure}

\subsection{Distribution shift robustness}
\label{sec:results_shift}

We also evaluate robustness to synthetic distribution shifts: Gaussian noise
and blur at two severity levels.

\begin{table}[H]
\centering
\caption{Effect of synthetic distribution shift on segmentation and
uncertainty quality. Under mild shift, Dice drops sharply but uncertainty
AUROC degrades more gradually. Under strong shift, both degrade substantially.}
\label{tab:shift}
\small
\setlength{\tabcolsep}{4.6pt}
\resizebox{\textwidth}{!}{%
\begin{tabular}{lccccc}
\toprule
\textbf{Shift level} & \textbf{Dice} & \textbf{AUC} & \textbf{ECE} &
\textbf{Unc-AUROC} & \textbf{ERR (adaptive)} \\
\midrule
None          & 0.764 & 0.965 & 0.037 & 0.765 & 45.6\% \\
Mild          & 0.517 & 0.840 & 0.056 & 0.648 & 7.9\% \\
Strong        & 0.186 & 0.758 & 0.094 & 0.655 & 25.0\% \\
\bottomrule
\end{tabular}%
}
\end{table}

Under mild shift ($\sigma_{\text{noise}} = 0.02$, $3\times3$ Gaussian blur
with $\sigma = 0.5$), Dice drops to 0.517, a 32\% relative decline. Unc-AUROC
drops to 0.648, still above chance but substantially degraded. Under strong
shift ($\sigma_{\text{noise}} = 0.10$, $7\times7$ blur, 4$\times$ resolution
reduction), Dice collapses to 0.186 (the model can barely segment), and
Unc-AUROC is 0.655.

The asymmetry is notable: even when the model fails catastrophically (Dice
0.186), uncertainty maintains above-random discriminative power (AUROC 0.655
vs.\ 0.5 random). However, the deferral error reduction under strong shift
is limited, because most pixels are wrong and the uncertainty signal, while
above random, is not sharp enough to efficiently separate the remaining correct
predictions.

This suggests a practical policy: under domain shift, monitor both average
uncertainty magnitude and uncertainty AUROC. When AUROC drops below a
threshold (e.g., 0.70), the system should flag the entire image for manual
review rather than relying on pixel-level deferral.
